% ==============================================================
% Section 2: Related Work
% Author: Hùng
% File:   paper/sections/02_related.tex
% ==============================================================

\section{Related Work}\label{sec:related}

The assessment of bug report quality has been studied from multiple
angles over the past decade.  We organize related work into three
themes: (1)~traditional machine-learning approaches for bug-report
classification, (2)~recent LLM-based methods for improving or
evaluating bug reports, and (3)~studies that employ inter-rater
agreement metrics such as Cohen's Kappa.  For each theme we highlight
what has been achieved and what remains unaddressed, ultimately
motivating the research gap that this study fills.

% ------------------------------------------------------------------
\subsection{Bug-Report Classification with Traditional ML}
\label{sec:related:ml}

Early efforts focused on distinguishing \emph{valid} from
\emph{invalid} bug reports using conventional classifiers.
Fan~et~al.~\cite{fan2020chaff} trained Random Forest and SVM models
on handcrafted features extracted from five Bugzilla projects,
achieving an average AUC of 0.86 for validity classification.
However, their approach treats a bug report as a monolithic entity
and does not evaluate individual quality dimensions such as
\emph{Observed Behavior} (OB), \emph{Expected Behavior} (EB), or
\emph{Steps to Reproduce} (S2R).  No inter-rater agreement was
measured.

Song and Chaparro~\cite{song2020bee} proposed \textsc{BEE}, an SVM-based
tool that detects OB, EB, and S2R sentences within bug reports,
reporting F1 scores of 0.93 (OB), 0.73 (EB), and 0.77 (S2R).
Although BEE explicitly targets the three reproducibility
components, the evaluation relied solely on automated metrics;
the authors acknowledged the need for user studies and did not
measure Cohen's Kappa against developer judgments.

These studies demonstrate that the problem of assessing bug-report
quality has long been recognized, yet the evaluation paradigm
remained limited to automated classification metrics without
measuring agreement with human reviewers.

% ------------------------------------------------------------------
\subsection{LLM-Based Bug-Report Enhancement and Evaluation}
\label{sec:related:llm}

With the advent of large language models, researchers began
leveraging their text-understanding capabilities for bug-report
tasks.  Bo~et~al.~\cite{bo2024chatbr} combined a BERT classifier
with ChatGPT (GPT-3.5) to both assess and improve bug reports from
Eclipse OSS projects.  Their tool, \textsc{ChatBR}, improved report
completeness substantially---OB by 13.0\%, EB by 84.2\%, and S2R by
41.0\%---but evaluated improvements exclusively through automated
semantic-similarity metrics, without any form of human agreement
measurement.

Acharya and Ginde~\cite{acharya2025enhance} fine-tuned Qwen~2.5-7B
and Llama~3.2-3B with LoRA on 3,903 Bugzilla reports, achieving a
CTQRS score of 77\% and an SBERT similarity of 0.82.  While these
results are competitive with GPT-4o few-shot baselines, the authors
themselves identified the absence of human evaluation as a key
limitation and called for ``human-judged evaluations'' in future
work.

Both studies confirm the potential of LLMs for processing bug
reports on general open-source repositories, yet neither measures
inter-rater agreement (e.g., Cohen's Kappa) between an LLM assessor
and a developer reviewer.

% ------------------------------------------------------------------
\subsection{Inter-Rater Agreement in Bug-Report Quality Assessment}
\label{sec:related:kappa}

A smaller body of work has incorporated Cohen's Kappa ($\kappa$) to
quantify annotation reliability, but with important scope
restrictions.

Mahmud~et~al.~\cite{mahmud2025astrobr} developed \textsc{AstroBR},
which pairs GPT-4 with dynamic app UI analysis to evaluate
Steps-to-Reproduce quality.  The study reports strong human--human
Kappa values ($\kappa = 0.91$ for sentence identification and
$\kappa = 0.88$ for S2R mapping) when constructing ground truth.
However, AstroBR is restricted to \emph{Android applications only}
and evaluates only the S2R dimension; the authors explicitly
acknowledge that OB and EB assessment is deferred to future work.

Akyol~et~al.~\cite{akyol2026improbr} presented \textsc{ImproBR},
which uses DistilBERT combined with GPT-4o~mini and
Retrieval-Augmented Generation to improve Minecraft (Mojira) bug
reports.  Notably, ImproBR reports inter-rater Kappa across all
three quality dimensions: $\kappa = 0.642$ (OB), $\kappa = 0.801$
(EB), and $\kappa = 0.663$ (S2R), with a pooled Kappa of 0.62.
Despite covering all three dimensions, the agreement is measured
\emph{between human annotators only}, not between the LLM and a
developer, and the study is confined to the Mojira platform.

Saha~et~al.~\cite{saha2026bugscribe} reported the highest Kappa in
our evidence set ($\kappa = 0.98$ for S2R on the development set),
but this value again reflects human--human agreement during ground
truth construction.  BugScribe is also limited to Android
applications and does not measure LLM-versus-developer Kappa.

% ------------------------------------------------------------------
\subsection{Research Gap}
\label{sec:related:gap}

Table~\ref{tab:related_summary} summarizes the positioning of our
study relative to prior work across four key dimensions.

\begin{table}[htbp]
\centering
\caption{Comparison of related studies along key dimensions.
         $\checkmark$ = addressed, $\times$ = not addressed.}
\label{tab:related_summary}
\small
\begin{tabular}{lcccc}
\hline
\textbf{Study} &
  \textbf{OB/EB/S2R} &
  \textbf{LLM} &
  \textbf{Kappa} &
  \textbf{Gen.\ OSS} \\
\hline
Fan~\cite{fan2020chaff}             & $\times$    & $\times$ & $\times$ & $\checkmark$ \\
Song~\cite{song2020bee}             & $\checkmark$ & $\times$ & $\times$ & $\checkmark$ \\
Bo~\cite{bo2024chatbr}              & $\checkmark$ & $\checkmark$ & $\times$ & $\checkmark$ \\
Acharya~\cite{acharya2025enhance}   & $\times$    & $\checkmark$ & $\times$ & $\checkmark$ \\
Mahmud~\cite{mahmud2025astrobr}     & S2R only     & $\checkmark$ & H--H     & $\times$ \\
Akyol~\cite{akyol2026improbr}       & $\checkmark$ & $\checkmark$ & H--H     & $\times$ \\
Saha~\cite{saha2026bugscribe}       & $\checkmark$ & $\checkmark$ & H--H     & $\times$ \\
\hline
\textbf{This study}                 & $\checkmark$ & $\checkmark$ & \textbf{LLM-Dev} & $\checkmark$ \\
\hline
\end{tabular}
\\[4pt]
\raggedright\footnotesize H--H = Human--Human Kappa only.
\end{table}

As the table shows, although prior work has made significant progress
in individual dimensions, \textbf{no existing study measures Cohen's
Kappa between an LLM assessor and a developer reviewer for the
combined OB\,+\,EB\,+\,S2R quality of bug reports on general GitHub
open-source projects}.  Studies that do employ Kappa (AstroBR,
ImproBR, BugScribe) are limited to platform-specific contexts
(Android, Mojira) and measure only human--human agreement during
ground truth construction.  Studies conducted on general OSS
repositories (ChatBR, Acharya~\&~Ginde) rely exclusively on
automated metrics without any human agreement evaluation.

Our research addresses this gap by directly measuring LLM-versus-developer
Cohen's Kappa for a combined assessment of OB, EB, and S2R on
GitHub open-source bug reports.
